\subsection{Objetivos Especificos}

\begin{itemize}
    \item Diseñar la arquitectura del sistema y el modelo de datos geoespacial que permita almacenar y estructurar información de siniestros viales.
    \item Desarrollar los módulos de análisis estadístico para calcular índices de accidentalidad por comuna, franja horaria, tipo de vehículo y gravedad.
    \item Integrar un tablero de monitoreo y diagnóstico que permitan visualizar zonas críticas y tendencias temporales.
    \item Validar el desempeño, seguridad y escalabilidad de la plataforma mediante pruebas técnicas y métricas de calidad de software.
\end{itemize}

\subsection{Componentes complementarios de la plataforma}

Los objetivos anteriores delimitan el alcance académico del trabajo. En la implementación de ViaData — Medellín se incorporaron, además, módulos de apoyo que facilitan la consulta y la operación del sistema sin ampliar la formulación original de objetivos.

El mapa analítico complementa el tablero de monitoreo con visualización territorial directa: puntos de incidentes, capas de densidad, coroplética por comuna o barrio y filtros geoespaciales compartidos con el resto de la plataforma.

Para perfiles autenticados con rol de analista o administrador, el sistema incluye reportes imprimibles que consolidan en formato exportable los resultados del tablero, del mapa y del módulo de predicciones. Asimismo, un asistente conversacional permite formular consultas en lenguaje natural invocando las mismas rutas de la API del backend. Por último, la gestión de usuarios y roles ---ciudadano, analista, administrador y autoridad reservada--- regula el acceso al módulo predictivo, a los reportes y a la administración de cuentas, en coherencia con el objetivo específico de validación en materia de seguridad.

El cumplimiento detallado de los objetivos y de estos componentes se presenta en las secciones de resultados y conclusiones.
